% Kompilacja: najpierw suplement (pdflatex suplement x2), potem
% pdflatex main && bibtex main && pdflatex main && pdflatex main
% Dane liczbowe (definicje.tex), tabele (tabele/*.csv) i rysunki (rysunki/) generuje
% analiza_odpowiedzi_naukowa.py; tekst (main.tex, sekcja-*.tex) jest edytowany ręcznie.
\documentclass[11pt]{article}
\usepackage[utf8]{inputenc}
\usepackage[T1]{fontenc}
\usepackage[polish]{babel}
\usepackage[margin=2.5cm]{geometry}
\usepackage{booktabs,graphicx,xcolor,amsmath}
\usepackage{xr-hyper}% odwołania do tabel i rycin S1, S2... z suplement.aux
\usepackage[hidelinks]{hyperref}
\externaldocument{suplement}
\usepackage{doi}
% BIBLIOGRAFIA:
\usepackage[numbers,sort&compress]{natbib}
\bibliographystyle{plainnat}

\extrafloats{100}
\newcommand{\fillin}[1]{\textcolor{red}{[#1]}}
% Wygenerowano automatycznie przez analiza_odpowiedzi_naukowa.py -- plik jest NADPISYWANY, nie edytuj ręcznie.
% Wstaw w preambule: % Wygenerowano automatycznie przez analiza_odpowiedzi_naukowa.py -- plik jest NADPISYWANY, nie edytuj ręcznie.
% Wstaw w preambule: % Wygenerowano automatycznie przez analiza_odpowiedzi_naukowa.py -- plik jest NADPISYWANY, nie edytuj ręcznie.
% Wstaw w preambule: \input{definicje.tex}
%
% Liczby w notacji polskiej: przecinek dziesiętny {,} (działa w tekście i w trybie matematycznym),
% odstęp tysięcy \, (od 10 000), minus \ensuremath{-}, przedziały [a; b].
% Makra wartości p (\SignP, \WilcoxonP, \P...) zawierają relację i są do trybu matematycznego: $p \SignP$.
%
% Tabele: \tabelacsv{<kolumny tabular>}{<nazwa pliku bez .csv>} wczytuje \TabDir<nazwa>.csv
% (csvsimple-l3; separator ';'; pierwszy wiersz = nagłówek; komórki w składni LaTeX; komórki ze średnikiem
% są ujęte w klamry; prefiks \midrule w pierwszej komórce = linia oddzielająca grupy wierszy).
% Przypis pod tabelą: \uwagatabeli{...}. Katalogi \TabDir i \FigDir można zdefiniować przed \input.
\RequirePackage{booktabs,csvsimple-l3}
\providecommand{\TabDir}{tabele/}
\providecommand{\FigDir}{rysunki/}
\providecommand{\fillin}[1]{\textbf{[#1]}}
\providecommand{\tabelacsv}[2]{\csvloop{file={\TabDir #2.csv}, separator=semicolon, head,
  column names detection=false, after head=\begin{tabular}{#1}\toprule\csvlinetotablerow\\\midrule,
  command=\csvlinetotablerow, late after line=\\, late after last line=\\\bottomrule\end{tabular}}}
\providecommand{\uwagatabeli}[1]{\par\vspace{2pt}\parbox{\linewidth}{\footnotesize #1}}
%
% Stałe liczbowe
\newcommand{\DaysOverThreePct}{32{,}0}
\newcommand{\DaysOverThreeMultiModPct}{84{,}9}
\newcommand{\DailyQOne}{2}
\newcommand{\DailyQThree}{4}
\newcommand{\NResponses}{57\,997}
\newcommand{\NEmployees}{180}
\newcommand{\NItems}{164}
\newcommand{\NModules}{9}
\newcommand{\NSequences}{15\,270}
\newcommand{\NRejected}{0}
\newcommand{\DateStart}{maja 2023}
\newcommand{\DateEnd}{czerwca 2026}
\newcommand{\AccOverall}{86{,}6}
\newcommand{\AccFirst}{82{,}6}
\newcommand{\AttemptsMedian}{3}
\newcommand{\AttemptsMax}{17}
\newcommand{\SeqRepeatedPct}{78{,}1}
\newcommand{\SeqMinThreePct}{59{,}8}
\newcommand{\GapMedian}{37}
\newcommand{\GapQOne}{18}
\newcommand{\GapQThree}{90}
\newcommand{\DailyMedian}{3}
\newcommand{\RespPerEmpMedian}{195}
\newcommand{\RespPerEmpQOne}{70}
\newcommand{\RespPerEmpQThree}{575}
\newcommand{\SpanMonthsMedian}{18}
\newcommand{\ActiveLastQuarterPct}{58{,}3}
\newcommand{\WeekdayPct}{100{,}0}
\newcommand{\ItemsPerEmpMedian}{86}
\newcommand{\ActiveDaysMedian}{84}
\newcommand{\OptionsMedian}{3}
\newcommand{\NExcluded}{212}
\newcommand{\NExcludedItems}{2}
\newcommand{\NExcludedItemResp}{348}
\newcommand{\NExcludedShortRep}{88}
\newcommand{\MinGapDays}{7}
\newcommand{\NShortRep}{88}
\newcommand{\NShortRepSameDay}{88}
\newcommand{\ShortRepPct}{0{,}15}
\newcommand{\NShortRepHour}{80}
\newcommand{\ExcludedModules}{mnk-tst}
\newcommand{\NExcludedPilot}{644}
\newcommand{\NExcludedPilotLater}{2233}
\newcommand{\NExcludedPilotTotal}{2877}
\newcommand{\NPilotSeq}{513}
\newcommand{\NPilotEmp}{72}
\newcommand{\PilotStart}{24.04.2023}
\newcommand{\PilotEnd}{28.04.2023}
\newcommand{\DaysUpToThreePct}{68{,}0}
\newcommand{\DaysUpToFourPct}{98{,}8}
\newcommand{\DailyMax}{7}
\newcommand{\HourPOne}{9:05}
\newcommand{\HourPNinetyNine}{15:57}
\newcommand{\HoursNineSixteenPct}{99{,}3}
\newcommand{\HoursTenFourteenPct}{62{,}6}
\newcommand{\ReaskAfterCorrect}{73{,}5}
\newcommand{\ReaskAfterIncorrect}{74{,}5}
\newcommand{\GapAfterCorrect}{37}
\newcommand{\GapAfterIncorrect}{39}
\newcommand{\AccAttOne}{82{,}6}
\newcommand{\AccAttOneCI}{[81{,}1; 84{,}1]}
\newcommand{\AccAttMax}{92{,}7}
\newcommand{\AccAttMaxCI}{[90{,}6; 94{,}8]}
\newcommand{\AttMax}{10}
\newcommand{\CohortK}{5}
\newcommand{\CohortN}{4890}
\newcommand{\CohortAccOne}{81{,}9}
\newcommand{\CohortAccK}{89{,}0}
\newcommand{\CohortAccOneCI}{[79{,}6; 84{,}1]}
\newcommand{\CohortAccKCI}{[87{,}3; 90{,}7]}
\newcommand{\PowerLawB}{0{,}29}
\newcommand{\PowerLawRsq}{0{,}97}
\newcommand{\ORMOneAttempt}{1{,}39}
\newcommand{\ORMOneAttemptCI}{[1{,}33; 1{,}45]}
\newcommand{\PMOneAttempt}{< 0{,}001}
\newcommand{\ORMTwoAttempt}{1{,}35}
\newcommand{\ORMTwoAttemptCI}{[1{,}25; 1{,}45]}
\newcommand{\PMTwoAttempt}{< 0{,}001}
\newcommand{\ORMTwoPrev}{3{,}89}
\newcommand{\ORMTwoPrevCI}{[3{,}51; 4{,}31]}
\newcommand{\PMTwoPrev}{< 0{,}001}
\newcommand{\ORMTwoGap}{0{,}88}
\newcommand{\ORMTwoGapCI}{[0{,}86; 0{,}90]}
\newcommand{\PMTwoGap}{< 0{,}001}
\newcommand{\ORMTwoTenure}{0{,}86}
\newcommand{\ORMTwoTenureCI}{[0{,}76; 0{,}98]}
\newcommand{\PMTwoTenure}{= 0{,}020}
\newcommand{\ORMThreeTenure}{0{,}88}
\newcommand{\ORMThreeTenureCI}{[0{,}76; 1{,}01]}
\newcommand{\PMThreeTenure}{= 0{,}067}
\newcommand{\MOneDropped}{292}
\newcommand{\ICCEmployee}{0{,}10}
\newcommand{\ICCItem}{0{,}24}
\newcommand{\SDEmployee}{0{,}71}
\newcommand{\SDItem}{1{,}10}
\newcommand{\RsqMarginal}{0{,}022}
\newcommand{\RsqConditional}{0{,}36}
\newcommand{\ORRange}{1{,}29--1{,}51}
\newcommand{\ORSensMain}{1{,}39}
\newcommand{\ORSensMainCI}{[1{,}33; 1{,}45]}
\newcommand{\ORSensCohort}{1{,}36}
\newcommand{\ORSensCohortCI}{[1{,}27; 1{,}46]}
\newcommand{\ORSensTrunc}{1{,}39}
\newcommand{\ORSensTruncCI}{[1{,}33; 1{,}45]}
\newcommand{\ORSensLenThree}{1{,}31}
\newcommand{\ORSensLenThreeCI}{[1{,}22; 1{,}41]}
\newcommand{\ORSensLenFive}{1{,}35}
\newcommand{\ORSensLenFiveCI}{[1{,}26; 1{,}46]}
\newcommand{\ORSensLenEight}{1{,}41}
\newcommand{\ORSensLenEightCI}{[1{,}32; 1{,}51]}
\newcommand{\ORSensActive}{1{,}39}
\newcommand{\ORSensActiveCI}{[1{,}33; 1{,}46]}
\newcommand{\ORSensTenure}{1{,}51}
\newcommand{\ORSensTenureCI}{[1{,}42; 1{,}61]}
\newcommand{\ORSensGee}{1{,}29}
\newcommand{\ORSensGeeCI}{[1{,}24; 1{,}35]}
\newcommand{\ORSensMixed}{1{,}38}
\newcommand{\ORSensMixedCI}{[1{,}34; 1{,}42]}
\newcommand{\NCorrections}{3382}
\newcommand{\NForgetting}{2719}
\newcommand{\EquilibriumPct}{88{,}8}
\newcommand{\RetainPct}{92{,}6}
\newcommand{\RetainCI}{[91{,}8; 93{,}4]}
\newcommand{\ForgetPct}{7{,}4}
\newcommand{\RecoverPct}{58{,}3}
\newcommand{\RecoverCI}{[55{,}5; 61{,}2]}
\newcommand{\PersistPct}{41{,}7}
\newcommand{\SameDistractorPct}{71{,}8}
\newcommand{\SameDistractorCI}{[69{,}4; 74{,}1]}
\newcommand{\SameDistractorN}{2420}
\newcommand{\SameDistractorExpected}{47{,}2}
\newcommand{\SameDistractorExpectedPop}{60{,}0}
\newcommand{\CorrRunOne}{64{,}4}
\newcommand{\CorrRunOneCI}{[62{,}2; 66{,}8]}
\newcommand{\CorrRunTwo}{48{,}1}
\newcommand{\CorrRunThree}{41{,}8}
\newcommand{\CorrRunMax}{34{,}2}
\newcommand{\CorrRunMaxCI}{[27{,}6; 40{,}9]}
\newcommand{\CorrRunMaxN}{295}
\newcommand{\RunMax}{4}
\newcommand{\EverCorrectedPct}{81{,}8}
\newcommand{\FirstWrongSeqN}{2068}
\newcommand{\SameDistrGapShort}{70{,}6}
\newcommand{\SameDistrGapMid}{73{,}2}
\newcommand{\SameDistrGapLong}{71{,}6}
\newcommand{\PersistTopItems}{14}
\newcommand{\PersistTopItemsShare}{39{,}6}
\newcommand{\PersistItems}{145}
\newcommand{\RetShortCorrect}{95{,}2}
\newcommand{\RetLongCorrect}{88{,}5}
\newcommand{\RetShortIncorrect}{60{,}9}
\newcommand{\RetLongIncorrect}{58{,}1}
\newcommand{\ForgetShort}{4{,}8}
\newcommand{\ForgetLong}{11{,}5}
\newcommand{\ForgetRunOne}{11{,}9}
\newcommand{\ForgetRunOneCI}{[10{,}7; 13{,}0]}
\newcommand{\ForgetRunTwo}{7{,}6}
\newcommand{\ForgetRunThree}{5{,}0}
\newcommand{\ForgetRunMax}{3{,}2}
\newcommand{\ForgetRunMaxCI}{[2{,}7; 3{,}7]}
\newcommand{\ForgetFreshShort}{7{,}7}
\newcommand{\ForgetFreshLong}{13{,}9}
\newcommand{\ForgetStableShort}{2{,}3}
\newcommand{\ForgetStableLong}{5{,}2}
\newcommand{\RelapsePct}{24{,}9}
\newcommand{\RelapseCI}{[22{,}9; 26{,}8]}
\newcommand{\RelapseN}{2457}
\newcommand{\RelapseSamePct}{73{,}6}
\newcommand{\RelapseSameCI}{[69{,}5; 77{,}7]}
\newcommand{\RelapseSameN}{611}
\newcommand{\ForgetRunK}{4}
\newcommand{\EmpN}{148}
\newcommand{\EmpImproved}{101}
\newcommand{\EmpWorsened}{18}
\newcommand{\EmpUnchanged}{29}
\newcommand{\EmpMeanChange}{4{,}1}
\newcommand{\EmpMeanChangeCI}{[2{,}7; 5{,}5]}
\newcommand{\SignP}{< 0{,}001}
\newcommand{\WilcoxonP}{< 0{,}001}
\newcommand{\RankBiserial}{0{,}77}
\newcommand{\EmpSharePct}{82{,}2}
\newcommand{\ItemDifficultyMedian}{85{,}5}
\newcommand{\TopDistractorMedian}{71{,}9}
\newcommand{\NDistractors}{349}
\newcommand{\NNonFunctional}{196}
\newcommand{\NonFunctionalPct}{56{,}2}
\newcommand{\ItemsModerate}{102}
\newcommand{\ItemsGoodDisc}{50}
\newcommand{\ItemsDiscTested}{161}
\newcommand{\ItemsImproving}{43}
\newcommand{\ItemsWorsening}{1}
\newcommand{\ItemsTestable}{164}
\newcommand{\ItemDiscMedian}{0{,}23}
\newcommand{\ItemsLowDisc}{65}
\newcommand{\ItemsEasy}{58}
\newcommand{\ItemsHard}{4}
\newcommand{\TopItems}{10}
\newcommand{\DemoN}{180}
\newcommand{\DemoMatched}{151}
\newcommand{\DemoMatchedPct}{83{,}9}
\newcommand{\DemoUnmatched}{29}
\newcommand{\DemoWomen}{121}
\newcommand{\DemoWomenPct}{80{,}1}
\newcommand{\DemoMen}{30}
\newcommand{\DemoMenPct}{19{,}9}
\newcommand{\DemoSexMissing}{29}
\newcommand{\DemoAgeN}{150}
\newcommand{\DemoAgeMissing}{30}
\newcommand{\AgeMedian}{47}
\newcommand{\AgeQOne}{41}
\newcommand{\AgeQThree}{56}
\newcommand{\AgeMin}{22}
\newcommand{\AgeMax}{67}
\newcommand{\AgeMean}{47{,}2}
\newcommand{\AgeSD}{10{,}6}
\newcommand{\AgeModalGroup}{50--59}
\newcommand{\AgeModalGroupPct}{34{,}0}
\newcommand{\FirstN}{15\,270}
\newcommand{\FirstAcc}{82{,}6}
\newcommand{\FirstAccCI}{[81{,}2; 84{,}1]}
\newcommand{\FirstAccMinGroup}{30--39}
\newcommand{\FirstAccMin}{74{,}9}
\newcommand{\FirstAccMaxGroup}{50--59}
\newcommand{\FirstAccMax}{84{,}3}
\newcommand{\FirstAgeOR}{1{,}16}
\newcommand{\FirstAgeORCI}{[1{,}03; 1{,}31]}
\newcommand{\FirstAgeP}{= 0{,}012}
\newcommand{\FirstAgeN}{13\,578}
\newcommand{\FirstAgeEmp}{150}
\newcommand{\FirstAgeRawOR}{1{,}12}
\newcommand{\FirstAgeRawORCI}{[1{,}01; 1{,}24]}
\newcommand{\FirstAgeRawP}{= 0{,}039}
\newcommand{\SeqChange}{7{,}2}
\newcommand{\SeqChangeCI}{[5{,}5; 8{,}9]}
\newcommand{\SeqChangeMinGroup}{60--69}
\newcommand{\SeqChangeMin}{4{,}0}
\newcommand{\SeqChangeMaxGroup}{40--49}
\newcommand{\SeqChangeMax}{8{,}8}
\newcommand{\SeqAgeOR}{0{,}97}
\newcommand{\SeqAgeORCI}{[0{,}90; 1{,}05]}
\newcommand{\SeqAgeP}{= 0{,}484}
\newcommand{\SeqAgeN}{23\,140}
\newcommand{\BootReps}{2000}
\newcommand{\Alpha}{0{,}05}
\newcommand{\MinItemN}{30}
\newcommand{\Seed}{2026}
\newcommand{\MaxAttemptPooled}{10}
\newcommand{\VerPython}{3.14.4}
\newcommand{\VerPandas}{3.0.3}
\newcommand{\VerNumpy}{2.5.0}
\newcommand{\VerScipy}{1.18.0}
\newcommand{\VerStatsmodels}{0.14.6}
\newcommand{\VerMatplotlib}{3.11.0}

%
% Liczby w notacji polskiej: przecinek dziesiętny {,} (działa w tekście i w trybie matematycznym),
% odstęp tysięcy \, (od 10 000), minus \ensuremath{-}, przedziały [a; b].
% Makra wartości p (\SignP, \WilcoxonP, \P...) zawierają relację i są do trybu matematycznego: $p \SignP$.
%
% Tabele: \tabelacsv{<kolumny tabular>}{<nazwa pliku bez .csv>} wczytuje \TabDir<nazwa>.csv
% (csvsimple-l3; separator ';'; pierwszy wiersz = nagłówek; komórki w składni LaTeX; komórki ze średnikiem
% są ujęte w klamry; prefiks \midrule w pierwszej komórce = linia oddzielająca grupy wierszy).
% Przypis pod tabelą: \uwagatabeli{...}. Katalogi \TabDir i \FigDir można zdefiniować przed \input.
\RequirePackage{booktabs,csvsimple-l3}
\providecommand{\TabDir}{tabele/}
\providecommand{\FigDir}{rysunki/}
\providecommand{\fillin}[1]{\textbf{[#1]}}
\providecommand{\tabelacsv}[2]{\csvloop{file={\TabDir #2.csv}, separator=semicolon, head,
  column names detection=false, after head=\begin{tabular}{#1}\toprule\csvlinetotablerow\\\midrule,
  command=\csvlinetotablerow, late after line=\\, late after last line=\\\bottomrule\end{tabular}}}
\providecommand{\uwagatabeli}[1]{\par\vspace{2pt}\parbox{\linewidth}{\footnotesize #1}}
%
% Stałe liczbowe
\newcommand{\DaysOverThreePct}{32{,}0}
\newcommand{\DaysOverThreeMultiModPct}{84{,}9}
\newcommand{\DailyQOne}{2}
\newcommand{\DailyQThree}{4}
\newcommand{\NResponses}{57\,997}
\newcommand{\NEmployees}{180}
\newcommand{\NItems}{164}
\newcommand{\NModules}{9}
\newcommand{\NSequences}{15\,270}
\newcommand{\NRejected}{0}
\newcommand{\DateStart}{maja 2023}
\newcommand{\DateEnd}{czerwca 2026}
\newcommand{\AccOverall}{86{,}6}
\newcommand{\AccFirst}{82{,}6}
\newcommand{\AttemptsMedian}{3}
\newcommand{\AttemptsMax}{17}
\newcommand{\SeqRepeatedPct}{78{,}1}
\newcommand{\SeqMinThreePct}{59{,}8}
\newcommand{\GapMedian}{37}
\newcommand{\GapQOne}{18}
\newcommand{\GapQThree}{90}
\newcommand{\DailyMedian}{3}
\newcommand{\RespPerEmpMedian}{195}
\newcommand{\RespPerEmpQOne}{70}
\newcommand{\RespPerEmpQThree}{575}
\newcommand{\SpanMonthsMedian}{18}
\newcommand{\ActiveLastQuarterPct}{58{,}3}
\newcommand{\WeekdayPct}{100{,}0}
\newcommand{\ItemsPerEmpMedian}{86}
\newcommand{\ActiveDaysMedian}{84}
\newcommand{\OptionsMedian}{3}
\newcommand{\NExcluded}{212}
\newcommand{\NExcludedItems}{2}
\newcommand{\NExcludedItemResp}{348}
\newcommand{\NExcludedShortRep}{88}
\newcommand{\MinGapDays}{7}
\newcommand{\NShortRep}{88}
\newcommand{\NShortRepSameDay}{88}
\newcommand{\ShortRepPct}{0{,}15}
\newcommand{\NShortRepHour}{80}
\newcommand{\ExcludedModules}{mnk-tst}
\newcommand{\NExcludedPilot}{644}
\newcommand{\NExcludedPilotLater}{2233}
\newcommand{\NExcludedPilotTotal}{2877}
\newcommand{\NPilotSeq}{513}
\newcommand{\NPilotEmp}{72}
\newcommand{\PilotStart}{24.04.2023}
\newcommand{\PilotEnd}{28.04.2023}
\newcommand{\DaysUpToThreePct}{68{,}0}
\newcommand{\DaysUpToFourPct}{98{,}8}
\newcommand{\DailyMax}{7}
\newcommand{\HourPOne}{9:05}
\newcommand{\HourPNinetyNine}{15:57}
\newcommand{\HoursNineSixteenPct}{99{,}3}
\newcommand{\HoursTenFourteenPct}{62{,}6}
\newcommand{\ReaskAfterCorrect}{73{,}5}
\newcommand{\ReaskAfterIncorrect}{74{,}5}
\newcommand{\GapAfterCorrect}{37}
\newcommand{\GapAfterIncorrect}{39}
\newcommand{\AccAttOne}{82{,}6}
\newcommand{\AccAttOneCI}{[81{,}1; 84{,}1]}
\newcommand{\AccAttMax}{92{,}7}
\newcommand{\AccAttMaxCI}{[90{,}6; 94{,}8]}
\newcommand{\AttMax}{10}
\newcommand{\CohortK}{5}
\newcommand{\CohortN}{4890}
\newcommand{\CohortAccOne}{81{,}9}
\newcommand{\CohortAccK}{89{,}0}
\newcommand{\CohortAccOneCI}{[79{,}6; 84{,}1]}
\newcommand{\CohortAccKCI}{[87{,}3; 90{,}7]}
\newcommand{\PowerLawB}{0{,}29}
\newcommand{\PowerLawRsq}{0{,}97}
\newcommand{\ORMOneAttempt}{1{,}39}
\newcommand{\ORMOneAttemptCI}{[1{,}33; 1{,}45]}
\newcommand{\PMOneAttempt}{< 0{,}001}
\newcommand{\ORMTwoAttempt}{1{,}35}
\newcommand{\ORMTwoAttemptCI}{[1{,}25; 1{,}45]}
\newcommand{\PMTwoAttempt}{< 0{,}001}
\newcommand{\ORMTwoPrev}{3{,}89}
\newcommand{\ORMTwoPrevCI}{[3{,}51; 4{,}31]}
\newcommand{\PMTwoPrev}{< 0{,}001}
\newcommand{\ORMTwoGap}{0{,}88}
\newcommand{\ORMTwoGapCI}{[0{,}86; 0{,}90]}
\newcommand{\PMTwoGap}{< 0{,}001}
\newcommand{\ORMTwoTenure}{0{,}86}
\newcommand{\ORMTwoTenureCI}{[0{,}76; 0{,}98]}
\newcommand{\PMTwoTenure}{= 0{,}020}
\newcommand{\ORMThreeTenure}{0{,}88}
\newcommand{\ORMThreeTenureCI}{[0{,}76; 1{,}01]}
\newcommand{\PMThreeTenure}{= 0{,}067}
\newcommand{\MOneDropped}{292}
\newcommand{\ICCEmployee}{0{,}10}
\newcommand{\ICCItem}{0{,}24}
\newcommand{\SDEmployee}{0{,}71}
\newcommand{\SDItem}{1{,}10}
\newcommand{\RsqMarginal}{0{,}022}
\newcommand{\RsqConditional}{0{,}36}
\newcommand{\ORRange}{1{,}29--1{,}51}
\newcommand{\ORSensMain}{1{,}39}
\newcommand{\ORSensMainCI}{[1{,}33; 1{,}45]}
\newcommand{\ORSensCohort}{1{,}36}
\newcommand{\ORSensCohortCI}{[1{,}27; 1{,}46]}
\newcommand{\ORSensTrunc}{1{,}39}
\newcommand{\ORSensTruncCI}{[1{,}33; 1{,}45]}
\newcommand{\ORSensLenThree}{1{,}31}
\newcommand{\ORSensLenThreeCI}{[1{,}22; 1{,}41]}
\newcommand{\ORSensLenFive}{1{,}35}
\newcommand{\ORSensLenFiveCI}{[1{,}26; 1{,}46]}
\newcommand{\ORSensLenEight}{1{,}41}
\newcommand{\ORSensLenEightCI}{[1{,}32; 1{,}51]}
\newcommand{\ORSensActive}{1{,}39}
\newcommand{\ORSensActiveCI}{[1{,}33; 1{,}46]}
\newcommand{\ORSensTenure}{1{,}51}
\newcommand{\ORSensTenureCI}{[1{,}42; 1{,}61]}
\newcommand{\ORSensGee}{1{,}29}
\newcommand{\ORSensGeeCI}{[1{,}24; 1{,}35]}
\newcommand{\ORSensMixed}{1{,}38}
\newcommand{\ORSensMixedCI}{[1{,}34; 1{,}42]}
\newcommand{\NCorrections}{3382}
\newcommand{\NForgetting}{2719}
\newcommand{\EquilibriumPct}{88{,}8}
\newcommand{\RetainPct}{92{,}6}
\newcommand{\RetainCI}{[91{,}8; 93{,}4]}
\newcommand{\ForgetPct}{7{,}4}
\newcommand{\RecoverPct}{58{,}3}
\newcommand{\RecoverCI}{[55{,}5; 61{,}2]}
\newcommand{\PersistPct}{41{,}7}
\newcommand{\SameDistractorPct}{71{,}8}
\newcommand{\SameDistractorCI}{[69{,}4; 74{,}1]}
\newcommand{\SameDistractorN}{2420}
\newcommand{\SameDistractorExpected}{47{,}2}
\newcommand{\SameDistractorExpectedPop}{60{,}0}
\newcommand{\CorrRunOne}{64{,}4}
\newcommand{\CorrRunOneCI}{[62{,}2; 66{,}8]}
\newcommand{\CorrRunTwo}{48{,}1}
\newcommand{\CorrRunThree}{41{,}8}
\newcommand{\CorrRunMax}{34{,}2}
\newcommand{\CorrRunMaxCI}{[27{,}6; 40{,}9]}
\newcommand{\CorrRunMaxN}{295}
\newcommand{\RunMax}{4}
\newcommand{\EverCorrectedPct}{81{,}8}
\newcommand{\FirstWrongSeqN}{2068}
\newcommand{\SameDistrGapShort}{70{,}6}
\newcommand{\SameDistrGapMid}{73{,}2}
\newcommand{\SameDistrGapLong}{71{,}6}
\newcommand{\PersistTopItems}{14}
\newcommand{\PersistTopItemsShare}{39{,}6}
\newcommand{\PersistItems}{145}
\newcommand{\RetShortCorrect}{95{,}2}
\newcommand{\RetLongCorrect}{88{,}5}
\newcommand{\RetShortIncorrect}{60{,}9}
\newcommand{\RetLongIncorrect}{58{,}1}
\newcommand{\ForgetShort}{4{,}8}
\newcommand{\ForgetLong}{11{,}5}
\newcommand{\ForgetRunOne}{11{,}9}
\newcommand{\ForgetRunOneCI}{[10{,}7; 13{,}0]}
\newcommand{\ForgetRunTwo}{7{,}6}
\newcommand{\ForgetRunThree}{5{,}0}
\newcommand{\ForgetRunMax}{3{,}2}
\newcommand{\ForgetRunMaxCI}{[2{,}7; 3{,}7]}
\newcommand{\ForgetFreshShort}{7{,}7}
\newcommand{\ForgetFreshLong}{13{,}9}
\newcommand{\ForgetStableShort}{2{,}3}
\newcommand{\ForgetStableLong}{5{,}2}
\newcommand{\RelapsePct}{24{,}9}
\newcommand{\RelapseCI}{[22{,}9; 26{,}8]}
\newcommand{\RelapseN}{2457}
\newcommand{\RelapseSamePct}{73{,}6}
\newcommand{\RelapseSameCI}{[69{,}5; 77{,}7]}
\newcommand{\RelapseSameN}{611}
\newcommand{\ForgetRunK}{4}
\newcommand{\EmpN}{148}
\newcommand{\EmpImproved}{101}
\newcommand{\EmpWorsened}{18}
\newcommand{\EmpUnchanged}{29}
\newcommand{\EmpMeanChange}{4{,}1}
\newcommand{\EmpMeanChangeCI}{[2{,}7; 5{,}5]}
\newcommand{\SignP}{< 0{,}001}
\newcommand{\WilcoxonP}{< 0{,}001}
\newcommand{\RankBiserial}{0{,}77}
\newcommand{\EmpSharePct}{82{,}2}
\newcommand{\ItemDifficultyMedian}{85{,}5}
\newcommand{\TopDistractorMedian}{71{,}9}
\newcommand{\NDistractors}{349}
\newcommand{\NNonFunctional}{196}
\newcommand{\NonFunctionalPct}{56{,}2}
\newcommand{\ItemsModerate}{102}
\newcommand{\ItemsGoodDisc}{50}
\newcommand{\ItemsDiscTested}{161}
\newcommand{\ItemsImproving}{43}
\newcommand{\ItemsWorsening}{1}
\newcommand{\ItemsTestable}{164}
\newcommand{\ItemDiscMedian}{0{,}23}
\newcommand{\ItemsLowDisc}{65}
\newcommand{\ItemsEasy}{58}
\newcommand{\ItemsHard}{4}
\newcommand{\TopItems}{10}
\newcommand{\DemoN}{180}
\newcommand{\DemoMatched}{151}
\newcommand{\DemoMatchedPct}{83{,}9}
\newcommand{\DemoUnmatched}{29}
\newcommand{\DemoWomen}{121}
\newcommand{\DemoWomenPct}{80{,}1}
\newcommand{\DemoMen}{30}
\newcommand{\DemoMenPct}{19{,}9}
\newcommand{\DemoSexMissing}{29}
\newcommand{\DemoAgeN}{150}
\newcommand{\DemoAgeMissing}{30}
\newcommand{\AgeMedian}{47}
\newcommand{\AgeQOne}{41}
\newcommand{\AgeQThree}{56}
\newcommand{\AgeMin}{22}
\newcommand{\AgeMax}{67}
\newcommand{\AgeMean}{47{,}2}
\newcommand{\AgeSD}{10{,}6}
\newcommand{\AgeModalGroup}{50--59}
\newcommand{\AgeModalGroupPct}{34{,}0}
\newcommand{\FirstN}{15\,270}
\newcommand{\FirstAcc}{82{,}6}
\newcommand{\FirstAccCI}{[81{,}2; 84{,}1]}
\newcommand{\FirstAccMinGroup}{30--39}
\newcommand{\FirstAccMin}{74{,}9}
\newcommand{\FirstAccMaxGroup}{50--59}
\newcommand{\FirstAccMax}{84{,}3}
\newcommand{\FirstAgeOR}{1{,}16}
\newcommand{\FirstAgeORCI}{[1{,}03; 1{,}31]}
\newcommand{\FirstAgeP}{= 0{,}012}
\newcommand{\FirstAgeN}{13\,578}
\newcommand{\FirstAgeEmp}{150}
\newcommand{\FirstAgeRawOR}{1{,}12}
\newcommand{\FirstAgeRawORCI}{[1{,}01; 1{,}24]}
\newcommand{\FirstAgeRawP}{= 0{,}039}
\newcommand{\SeqChange}{7{,}2}
\newcommand{\SeqChangeCI}{[5{,}5; 8{,}9]}
\newcommand{\SeqChangeMinGroup}{60--69}
\newcommand{\SeqChangeMin}{4{,}0}
\newcommand{\SeqChangeMaxGroup}{40--49}
\newcommand{\SeqChangeMax}{8{,}8}
\newcommand{\SeqAgeOR}{0{,}97}
\newcommand{\SeqAgeORCI}{[0{,}90; 1{,}05]}
\newcommand{\SeqAgeP}{= 0{,}484}
\newcommand{\SeqAgeN}{23\,140}
\newcommand{\BootReps}{2000}
\newcommand{\Alpha}{0{,}05}
\newcommand{\MinItemN}{30}
\newcommand{\Seed}{2026}
\newcommand{\MaxAttemptPooled}{10}
\newcommand{\VerPython}{3.14.4}
\newcommand{\VerPandas}{3.0.3}
\newcommand{\VerNumpy}{2.5.0}
\newcommand{\VerScipy}{1.18.0}
\newcommand{\VerStatsmodels}{0.14.6}
\newcommand{\VerMatplotlib}{3.11.0}

%
% Liczby w notacji polskiej: przecinek dziesiętny {,} (działa w tekście i w trybie matematycznym),
% odstęp tysięcy \, (od 10 000), minus \ensuremath{-}, przedziały [a; b].
% Makra wartości p (\SignP, \WilcoxonP, \P...) zawierają relację i są do trybu matematycznego: $p \SignP$.
%
% Tabele: \tabelacsv{<kolumny tabular>}{<nazwa pliku bez .csv>} wczytuje \TabDir<nazwa>.csv
% (csvsimple-l3; separator ';'; pierwszy wiersz = nagłówek; komórki w składni LaTeX; komórki ze średnikiem
% są ujęte w klamry; prefiks \midrule w pierwszej komórce = linia oddzielająca grupy wierszy).
% Przypis pod tabelą: \uwagatabeli{...}. Katalogi \TabDir i \FigDir można zdefiniować przed \input.
\RequirePackage{booktabs,csvsimple-l3}
\providecommand{\TabDir}{tabele/}
\providecommand{\FigDir}{rysunki/}
\providecommand{\fillin}[1]{\textbf{[#1]}}
\providecommand{\tabelacsv}[2]{\csvloop{file={\TabDir #2.csv}, separator=semicolon, head,
  column names detection=false, after head=\begin{tabular}{#1}\toprule\csvlinetotablerow\\\midrule,
  command=\csvlinetotablerow, late after line=\\, late after last line=\\\bottomrule\end{tabular}}}
\providecommand{\uwagatabeli}[1]{\par\vspace{2pt}\parbox{\linewidth}{\footnotesize #1}}
%
% Stałe liczbowe
\newcommand{\DaysOverThreePct}{32{,}0}
\newcommand{\DaysOverThreeMultiModPct}{84{,}9}
\newcommand{\DailyQOne}{2}
\newcommand{\DailyQThree}{4}
\newcommand{\NResponses}{57\,997}
\newcommand{\NEmployees}{180}
\newcommand{\NItems}{164}
\newcommand{\NModules}{9}
\newcommand{\NSequences}{15\,270}
\newcommand{\NRejected}{0}
\newcommand{\DateStart}{maja 2023}
\newcommand{\DateEnd}{czerwca 2026}
\newcommand{\AccOverall}{86{,}6}
\newcommand{\AccFirst}{82{,}6}
\newcommand{\AttemptsMedian}{3}
\newcommand{\AttemptsMax}{17}
\newcommand{\SeqRepeatedPct}{78{,}1}
\newcommand{\SeqMinThreePct}{59{,}8}
\newcommand{\GapMedian}{37}
\newcommand{\GapQOne}{18}
\newcommand{\GapQThree}{90}
\newcommand{\DailyMedian}{3}
\newcommand{\RespPerEmpMedian}{195}
\newcommand{\RespPerEmpQOne}{70}
\newcommand{\RespPerEmpQThree}{575}
\newcommand{\SpanMonthsMedian}{18}
\newcommand{\ActiveLastQuarterPct}{58{,}3}
\newcommand{\WeekdayPct}{100{,}0}
\newcommand{\ItemsPerEmpMedian}{86}
\newcommand{\ActiveDaysMedian}{84}
\newcommand{\OptionsMedian}{3}
\newcommand{\NExcluded}{212}
\newcommand{\NExcludedItems}{2}
\newcommand{\NExcludedItemResp}{348}
\newcommand{\NExcludedShortRep}{88}
\newcommand{\MinGapDays}{7}
\newcommand{\NShortRep}{88}
\newcommand{\NShortRepSameDay}{88}
\newcommand{\ShortRepPct}{0{,}15}
\newcommand{\NShortRepHour}{80}
\newcommand{\ExcludedModules}{mnk-tst}
\newcommand{\NExcludedPilot}{644}
\newcommand{\NExcludedPilotLater}{2233}
\newcommand{\NExcludedPilotTotal}{2877}
\newcommand{\NPilotSeq}{513}
\newcommand{\NPilotEmp}{72}
\newcommand{\PilotStart}{24.04.2023}
\newcommand{\PilotEnd}{28.04.2023}
\newcommand{\DaysUpToThreePct}{68{,}0}
\newcommand{\DaysUpToFourPct}{98{,}8}
\newcommand{\DailyMax}{7}
\newcommand{\HourPOne}{9:05}
\newcommand{\HourPNinetyNine}{15:57}
\newcommand{\HoursNineSixteenPct}{99{,}3}
\newcommand{\HoursTenFourteenPct}{62{,}6}
\newcommand{\ReaskAfterCorrect}{73{,}5}
\newcommand{\ReaskAfterIncorrect}{74{,}5}
\newcommand{\GapAfterCorrect}{37}
\newcommand{\GapAfterIncorrect}{39}
\newcommand{\AccAttOne}{82{,}6}
\newcommand{\AccAttOneCI}{[81{,}1; 84{,}1]}
\newcommand{\AccAttMax}{92{,}7}
\newcommand{\AccAttMaxCI}{[90{,}6; 94{,}8]}
\newcommand{\AttMax}{10}
\newcommand{\CohortK}{5}
\newcommand{\CohortN}{4890}
\newcommand{\CohortAccOne}{81{,}9}
\newcommand{\CohortAccK}{89{,}0}
\newcommand{\CohortAccOneCI}{[79{,}6; 84{,}1]}
\newcommand{\CohortAccKCI}{[87{,}3; 90{,}7]}
\newcommand{\PowerLawB}{0{,}29}
\newcommand{\PowerLawRsq}{0{,}97}
\newcommand{\ORMOneAttempt}{1{,}39}
\newcommand{\ORMOneAttemptCI}{[1{,}33; 1{,}45]}
\newcommand{\PMOneAttempt}{< 0{,}001}
\newcommand{\ORMTwoAttempt}{1{,}35}
\newcommand{\ORMTwoAttemptCI}{[1{,}25; 1{,}45]}
\newcommand{\PMTwoAttempt}{< 0{,}001}
\newcommand{\ORMTwoPrev}{3{,}89}
\newcommand{\ORMTwoPrevCI}{[3{,}51; 4{,}31]}
\newcommand{\PMTwoPrev}{< 0{,}001}
\newcommand{\ORMTwoGap}{0{,}88}
\newcommand{\ORMTwoGapCI}{[0{,}86; 0{,}90]}
\newcommand{\PMTwoGap}{< 0{,}001}
\newcommand{\ORMTwoTenure}{0{,}86}
\newcommand{\ORMTwoTenureCI}{[0{,}76; 0{,}98]}
\newcommand{\PMTwoTenure}{= 0{,}020}
\newcommand{\ORMThreeTenure}{0{,}88}
\newcommand{\ORMThreeTenureCI}{[0{,}76; 1{,}01]}
\newcommand{\PMThreeTenure}{= 0{,}067}
\newcommand{\MOneDropped}{292}
\newcommand{\ICCEmployee}{0{,}10}
\newcommand{\ICCItem}{0{,}24}
\newcommand{\SDEmployee}{0{,}71}
\newcommand{\SDItem}{1{,}10}
\newcommand{\RsqMarginal}{0{,}022}
\newcommand{\RsqConditional}{0{,}36}
\newcommand{\ORRange}{1{,}29--1{,}51}
\newcommand{\ORSensMain}{1{,}39}
\newcommand{\ORSensMainCI}{[1{,}33; 1{,}45]}
\newcommand{\ORSensCohort}{1{,}36}
\newcommand{\ORSensCohortCI}{[1{,}27; 1{,}46]}
\newcommand{\ORSensTrunc}{1{,}39}
\newcommand{\ORSensTruncCI}{[1{,}33; 1{,}45]}
\newcommand{\ORSensLenThree}{1{,}31}
\newcommand{\ORSensLenThreeCI}{[1{,}22; 1{,}41]}
\newcommand{\ORSensLenFive}{1{,}35}
\newcommand{\ORSensLenFiveCI}{[1{,}26; 1{,}46]}
\newcommand{\ORSensLenEight}{1{,}41}
\newcommand{\ORSensLenEightCI}{[1{,}32; 1{,}51]}
\newcommand{\ORSensActive}{1{,}39}
\newcommand{\ORSensActiveCI}{[1{,}33; 1{,}46]}
\newcommand{\ORSensTenure}{1{,}51}
\newcommand{\ORSensTenureCI}{[1{,}42; 1{,}61]}
\newcommand{\ORSensGee}{1{,}29}
\newcommand{\ORSensGeeCI}{[1{,}24; 1{,}35]}
\newcommand{\ORSensMixed}{1{,}38}
\newcommand{\ORSensMixedCI}{[1{,}34; 1{,}42]}
\newcommand{\NCorrections}{3382}
\newcommand{\NForgetting}{2719}
\newcommand{\EquilibriumPct}{88{,}8}
\newcommand{\RetainPct}{92{,}6}
\newcommand{\RetainCI}{[91{,}8; 93{,}4]}
\newcommand{\ForgetPct}{7{,}4}
\newcommand{\RecoverPct}{58{,}3}
\newcommand{\RecoverCI}{[55{,}5; 61{,}2]}
\newcommand{\PersistPct}{41{,}7}
\newcommand{\SameDistractorPct}{71{,}8}
\newcommand{\SameDistractorCI}{[69{,}4; 74{,}1]}
\newcommand{\SameDistractorN}{2420}
\newcommand{\SameDistractorExpected}{47{,}2}
\newcommand{\SameDistractorExpectedPop}{60{,}0}
\newcommand{\CorrRunOne}{64{,}4}
\newcommand{\CorrRunOneCI}{[62{,}2; 66{,}8]}
\newcommand{\CorrRunTwo}{48{,}1}
\newcommand{\CorrRunThree}{41{,}8}
\newcommand{\CorrRunMax}{34{,}2}
\newcommand{\CorrRunMaxCI}{[27{,}6; 40{,}9]}
\newcommand{\CorrRunMaxN}{295}
\newcommand{\RunMax}{4}
\newcommand{\EverCorrectedPct}{81{,}8}
\newcommand{\FirstWrongSeqN}{2068}
\newcommand{\SameDistrGapShort}{70{,}6}
\newcommand{\SameDistrGapMid}{73{,}2}
\newcommand{\SameDistrGapLong}{71{,}6}
\newcommand{\PersistTopItems}{14}
\newcommand{\PersistTopItemsShare}{39{,}6}
\newcommand{\PersistItems}{145}
\newcommand{\RetShortCorrect}{95{,}2}
\newcommand{\RetLongCorrect}{88{,}5}
\newcommand{\RetShortIncorrect}{60{,}9}
\newcommand{\RetLongIncorrect}{58{,}1}
\newcommand{\ForgetShort}{4{,}8}
\newcommand{\ForgetLong}{11{,}5}
\newcommand{\ForgetRunOne}{11{,}9}
\newcommand{\ForgetRunOneCI}{[10{,}7; 13{,}0]}
\newcommand{\ForgetRunTwo}{7{,}6}
\newcommand{\ForgetRunThree}{5{,}0}
\newcommand{\ForgetRunMax}{3{,}2}
\newcommand{\ForgetRunMaxCI}{[2{,}7; 3{,}7]}
\newcommand{\ForgetFreshShort}{7{,}7}
\newcommand{\ForgetFreshLong}{13{,}9}
\newcommand{\ForgetStableShort}{2{,}3}
\newcommand{\ForgetStableLong}{5{,}2}
\newcommand{\RelapsePct}{24{,}9}
\newcommand{\RelapseCI}{[22{,}9; 26{,}8]}
\newcommand{\RelapseN}{2457}
\newcommand{\RelapseSamePct}{73{,}6}
\newcommand{\RelapseSameCI}{[69{,}5; 77{,}7]}
\newcommand{\RelapseSameN}{611}
\newcommand{\ForgetRunK}{4}
\newcommand{\EmpN}{148}
\newcommand{\EmpImproved}{101}
\newcommand{\EmpWorsened}{18}
\newcommand{\EmpUnchanged}{29}
\newcommand{\EmpMeanChange}{4{,}1}
\newcommand{\EmpMeanChangeCI}{[2{,}7; 5{,}5]}
\newcommand{\SignP}{< 0{,}001}
\newcommand{\WilcoxonP}{< 0{,}001}
\newcommand{\RankBiserial}{0{,}77}
\newcommand{\EmpSharePct}{82{,}2}
\newcommand{\ItemDifficultyMedian}{85{,}5}
\newcommand{\TopDistractorMedian}{71{,}9}
\newcommand{\NDistractors}{349}
\newcommand{\NNonFunctional}{196}
\newcommand{\NonFunctionalPct}{56{,}2}
\newcommand{\ItemsModerate}{102}
\newcommand{\ItemsGoodDisc}{50}
\newcommand{\ItemsDiscTested}{161}
\newcommand{\ItemsImproving}{43}
\newcommand{\ItemsWorsening}{1}
\newcommand{\ItemsTestable}{164}
\newcommand{\ItemDiscMedian}{0{,}23}
\newcommand{\ItemsLowDisc}{65}
\newcommand{\ItemsEasy}{58}
\newcommand{\ItemsHard}{4}
\newcommand{\TopItems}{10}
\newcommand{\DemoN}{180}
\newcommand{\DemoMatched}{151}
\newcommand{\DemoMatchedPct}{83{,}9}
\newcommand{\DemoUnmatched}{29}
\newcommand{\DemoWomen}{121}
\newcommand{\DemoWomenPct}{80{,}1}
\newcommand{\DemoMen}{30}
\newcommand{\DemoMenPct}{19{,}9}
\newcommand{\DemoSexMissing}{29}
\newcommand{\DemoAgeN}{150}
\newcommand{\DemoAgeMissing}{30}
\newcommand{\AgeMedian}{47}
\newcommand{\AgeQOne}{41}
\newcommand{\AgeQThree}{56}
\newcommand{\AgeMin}{22}
\newcommand{\AgeMax}{67}
\newcommand{\AgeMean}{47{,}2}
\newcommand{\AgeSD}{10{,}6}
\newcommand{\AgeModalGroup}{50--59}
\newcommand{\AgeModalGroupPct}{34{,}0}
\newcommand{\FirstN}{15\,270}
\newcommand{\FirstAcc}{82{,}6}
\newcommand{\FirstAccCI}{[81{,}2; 84{,}1]}
\newcommand{\FirstAccMinGroup}{30--39}
\newcommand{\FirstAccMin}{74{,}9}
\newcommand{\FirstAccMaxGroup}{50--59}
\newcommand{\FirstAccMax}{84{,}3}
\newcommand{\FirstAgeOR}{1{,}16}
\newcommand{\FirstAgeORCI}{[1{,}03; 1{,}31]}
\newcommand{\FirstAgeP}{= 0{,}012}
\newcommand{\FirstAgeN}{13\,578}
\newcommand{\FirstAgeEmp}{150}
\newcommand{\FirstAgeRawOR}{1{,}12}
\newcommand{\FirstAgeRawORCI}{[1{,}01; 1{,}24]}
\newcommand{\FirstAgeRawP}{= 0{,}039}
\newcommand{\SeqChange}{7{,}2}
\newcommand{\SeqChangeCI}{[5{,}5; 8{,}9]}
\newcommand{\SeqChangeMinGroup}{60--69}
\newcommand{\SeqChangeMin}{4{,}0}
\newcommand{\SeqChangeMaxGroup}{40--49}
\newcommand{\SeqChangeMax}{8{,}8}
\newcommand{\SeqAgeOR}{0{,}97}
\newcommand{\SeqAgeORCI}{[0{,}90; 1{,}05]}
\newcommand{\SeqAgeP}{= 0{,}484}
\newcommand{\SeqAgeN}{23\,140}
\newcommand{\BootReps}{2000}
\newcommand{\Alpha}{0{,}05}
\newcommand{\MinItemN}{30}
\newcommand{\Seed}{2026}
\newcommand{\MaxAttemptPooled}{10}
\newcommand{\VerPython}{3.14.4}
\newcommand{\VerPandas}{3.0.3}
\newcommand{\VerNumpy}{2.5.0}
\newcommand{\VerScipy}{1.18.0}
\newcommand{\VerStatsmodels}{0.14.6}
\newcommand{\VerMatplotlib}{3.11.0}


\usepackage{authblk}
\renewcommand\Authand{ i }\renewcommand\Authands{ i }
\author[1,*]{Bartosz Lewandowski \small{(ORCID:0009-0004-1158-0693)}}
\author[2]{Monika Szerszeń \small{(ORCID: 0009-0009-6625-3979)}}
\author[3]{Magdalena Dowgalec \small{(ORCID:0009-0005-4469-2832)}}
\author[4]{Ewaryst Tkacz \small{(ORCID:0000-0003-2580-7954)}}
\author[5]{Maciej J. Wróbel \small{(ORCID:0000-0003-1865-2122)}}

\affil[1,4]{Department of Clinical Engineering, Academy of Silesia in Katowice}
\affil[3]{Powiślański Bank Spółdzielczy w Kwidzynie}
\affil[4]{Atende SA}
\affil[5]{Department of Otolaryngology and Oncological Laryngology Collegium Medicum NCU}

% Autor korespondencyjny
\affil[*]{Autor korespondencyjny: dr Bartosz Lewandowski, bartosz.lewandowski@akademiaslaska.pl}

\title{Powtarzane pytania, poprawność odpowiedzi i trwałość błędów: trzyletnia analiza mikrouczenia w organizacji}


\date{}
\begin{document}
\maketitle


\begin{abstract}
\noindent\textbf{Wprowadzenie.} Przypominanie rozłożone w czasie poprawia zapamiętywanie w warunkach
laboratoryjnych i edukacyjnych, jednak wyniki wieloletnich programów szkoleń w miejscu pracy opisano
w niewielu badaniach.
\textbf{Metody.} Przeanalizowano \NResponses{} odpowiedzi udzielonych przez \NEmployees{} pracowników
banku na \NItems{} pytań wielokrotnego wyboru w programie mikrouczenia opartym na quizach (od
\DateStart{} do \DateEnd). Pytania wracały do pracownika w stałym, nieadaptacyjnym odstępie (mediana
\GapMedian{} dni), a po każdej odpowiedzi wyświetlano odpowiedź poprawną. Poprawność modelowano
regresją logistyczną z efektami stałymi pracownika i pytania w trzech modelach: uczenia, retencji
i nowych pytań.
\textbf{Wyniki.} Odsetek poprawnych odpowiedzi wzrósł z \AccAttOne\% w pierwszej próbie do
\AccAttMax\% w próbie \AttMax{}. i kolejnych. Każde podwojenie numeru próby zwiększało szanse
poprawnej odpowiedzi \ORMOneAttempt-krotnie (95\% PU \ORMOneAttemptCI; zakres w analizach
wrażliwości \ORRange), a każde podwojenie odstępu od poprzedniej próby zmniejszało je (OR
\ORMTwoGap). Po powtórzonym błędzie pracownicy w \SameDistractorPct\% przypadków wybierali ten sam
dystraktor, wobec \SameDistractorExpectedPop\% oczekiwanych przy uwzględnieniu popularności
dystraktorów. Poprawność na nowo napotykanych pytaniach nie rosła wraz ze stażem (OR
\ORMThreeTenure{} na rok, 95\% PU \ORMThreeTenureCI). % SPRAWDŹ
\textbf{Wnioski.} Powtarzanemu przypominaniu towarzyszyła poprawa w powtarzanych pytaniach,
zapominanie między powtórkami i utrwalanie konkretnych błędów mimo informacji zwrotnej. Dane nie
rozstrzygają, czy poprawa dotyczy wiedzy, czy zapamiętania odpowiedzi. Adaptacyjny harmonogram
powtórek i informacja zwrotna odnosząca się do wybranego dystraktora mogłyby zwiększyć skuteczność
takich programów.

\medskip\noindent\textbf{Słowa kluczowe:} powtórki rozłożone w czasie; praktyka przypominania;
mikrouczenie; szkolenia pracownicze; analityka uczenia się; informacja zwrotna; zapominanie
\end{abstract}

% =============================================================================
% Wprowadzenie. Plik edytowany ręcznie -- skrypt analiza_odpowiedzi_naukowa.py go NIE nadpisuje.
% Liczby: definicje.tex; tabele: \TabDir (CSV); rysunki: \FigDir; literatura: references.bib.
% =============================================================================
\section{Wprowadzenie}

Przypominanie sobie informacji z pamięci wzmacnia ich późniejsze utrwalenie silniej niż ponowne
ich przeczytanie (efekt testowania), a nauka rozłożona w czasie jest zapamiętywana lepiej niż nauka
skumulowana (efekt rozłożenia w czasie) \citep{roediger2006,cepeda2006,cepeda2008}. W randomizowanych
badaniach z udziałem lekarzy w trakcie kształcenia tzw. edukacja rozłożona w czasie (ang.
\emph{spaced education}), w której krótkie pytania wielokrotnego wyboru przesyła się cyklicznie
przez tygodnie lub miesiące, poprawiała przyswajanie i utrwalanie wiedzy
\citep{kerfoot2007urology,kerfoot2007mededuc,kerfoot2009}.

W szkoleniach organizacyjnych stosuje się \emph{mikrouczenie} (ang. \emph{microlearning}), czyli
krótkie jednostki szkoleniowe, często w formie quizów. Dowody jego skuteczności pochodzą jednak głównie
z badań na małych próbach lub z tzw. szarej literatury \citep{taylor2022}, a ewaluacje obowiązkowych
szkoleń organizacyjnych, np.\ z zakresu świadomości bezpieczeństwa, często wykazują efekty niewielkie
albo zanikające z upływem czasu \citep{reinheimer2020,lain2022,ho2025}. Otwarte pozostaje więc
pytanie, czy wiedza sprawdzana w trakcie takich szkoleń jest nabywana i utrwalana.

Skuteczność szkoleń organizacyjnych ocenia się zwykle przez transfer, rozumiany jako zastosowanie
nabytej wiedzy i umiejętności w pracy \citep{baldwin1988,blume2010}. Warunkiem transferu jest jednak
wcześniejsze nabycie i utrwalenie wiedzy, a ten etap można zweryfikować na podstawie danych
gromadzonych przez platformy mikrouczenia. Platformy te zapisują każdą odpowiedź wraz z czasem
i wynikiem, a to samo pytanie wraca do tego samego pracownika wielokrotnie. Kolejne próby tworzą
zatem serię pomiarów tej samej wiedzy u tej samej osoby, co pozwala ocenić, czy poprawność rośnie,
czy wiedza utrzymuje się między powtórkami i czy błędy się utrwalają. Dane na poziomie pojedynczych
odpowiedzi wykorzystuje się w podobny sposób w modelach ucznia rozwijanych dla inteligentnych systemów
tutorskich i aplikacji do nauki języków \citep{corbett1994,pavlik2009,settles2016,tabibian2019,koedinger2023}.
W niniejszym badaniu podejście to zastosowano do walidacji wiedzy nabywanej w ramach mikrouczenia
w organizacji.

Badanie przeprowadzono w banku, w którym pracownicy odpowiadali na krótkie pytania wielokrotnego
wyboru dotyczące wymogów regulacyjnych, produktów bankowych, umiejętności miękkich oraz, w przypadku
kadry kierowniczej, zarządzania, wyświetlane przez aplikację na komputerze w trakcie dnia pracy.
Pytania wracały do pracownika w stałym, nieadaptacyjnym odstępie, a po każdej odpowiedzi wyświetlano
odpowiedź poprawną. Na podstawie analizy odpowiedzi z okresu trzech lat,  sformułowano trzy pytania badawcze:
(PB1) czy odsetek poprawnych odpowiedzi rośnie wraz z kolejnymi przypomnieniami tego samego pytania
przez tego samego pracownika;
(PB2) jak poprawność zależy od wyniku poprzedniej próby i od czasu, który od niej upłynął;
(PB3) czy powtórzone błędy są trwałe, co wskazywałoby raczej na utrwalone błędne przekonania niż na
zgadywanie, mimo informacji zwrotnej po każdej odpowiedzi.

% =============================================================================
% Materiał i metody. Plik edytowany ręcznie -- skrypt analiza_odpowiedzi_naukowa.py go NIE nadpisuje.
% Liczby: definicje.tex; tabele: \TabDir (CSV); rysunki: \FigDir; literatura: references.bib.
% Analizy uzupełniające (zmiana na poziomie pracowników, analiza pytań) opisano w suplement.tex.
% =============================================================================
\section{Materiał i metody}

\subsection{Środowisko i program szkoleniowy}
Program prowadzono w banku. Jego treść stanowiły pytania ogólne pogrupowane w cztery kategorie:
wymogi regulacyjne, produkty bankowe, umiejętności miękkie oraz moduł dla kadry kierowniczej.
Każde pytanie było pytaniem wielokrotnego wyboru z trzema do pięciu opcjami odpowiedzi, z których
dokładnie jedna była poprawna, i zawierało wyjaśnienie poprawnej odpowiedzi. Aplikacja szkoleniowa
uruchamiała się automatycznie na ekranach komputerów pracowników i wyświetlała od jednego do trzech
pytań dziennie w godzinach 10:00--14:00; po każdej odpowiedzi pokazywano poprawną odpowiedź wraz
z wyjaśnieniem. To samo pytanie mogło zostać ponownie zadane temu samemu pracownikowi nie wcześniej
niż siedem dni po poprzedniej odpowiedzi, a harmonogram powtórek nie zależał od tego, czy odpowiedź
była poprawna. Na czas przeglądu i aktualizacji pytań program był zawieszany.
O udziale w programie decydował Departament Digitalizacji i Procesów banku, który był również autorem
pytań. Pytania prezentowano w modułach tematycznych (mediana \OptionsMedian{} zaobserwowanych opcji
odpowiedzi na pytanie). W analizowanych danych pracownicy odpowiadali na medianę \DailyMedian{}
pytań w dniu aktywności, a to samo pytanie było ponownie zadawane temu samemu pracownikowi po
medianie \GapMedian{} dni (IQR \GapQOne--\GapQThree).

\subsection{Dane}
Dane przekazano w formie listy odpowiedzi zawierającej moduł, pytanie, wybraną opcję, kod
odpowiadającego, znacznik czasu oraz poprawność odpowiedzi (wartość logiczna). Wykluczono
\NExcludedItems{} pytania, w których zarejestrowano tylko jedną opcję odpowiedzi
(\NExcludedItemResp{} odpowiedzi), ponieważ nie niosą one informacji o błędach. Zbiór analityczny
obejmował \NResponses{} odpowiedzi (tab.~\ref{tab:zbior_danych}). Dane osób odpowiadających
przekazano w formie metadanych zawierających płeć pracownika oraz przedział wiekowy (co 10 lat) oraz medianę wieku.
% SPRAWDŹ: people.csv - nie bedzie dostarczany i brany pod uwage.

\fillin{Podstawa prawna przetwarzania danych spseudonimizowanych (wiek, płeć)}

\subsection{Definicje}
\emph{Sekwencja} to uporządkowana chronologicznie seria prób jednego pracownika dotyczących jednego
pytania; \emph{numer próby} to pozycja odpowiedzi w jej sekwencji (przy jednakowym czasie o kolejności
decydował identyfikator rekordu). \emph{Poprawność} to odsetek poprawnych odpowiedzi. \emph{Odstęp}
to czas w dniach od poprzedniej próby w tej samej sekwencji. \emph{Staż} to czas od pierwszej
odpowiedzi pracownika w programie.

\subsection{Analiza opisowa}
Odsetki poprawnych odpowiedzi podano z 95-procentowymi przedziałami ufności (PU) z bootstrapu
klastrowego, w którym losowano pracowników ze zwracaniem (\BootReps{} replikacji); procedura ta
uwzględnia zależność odpowiedzi udzielanych przez tę samą osobę. Ponieważ osoby kontynuujące do
kolejnych prób mogą różnić się od tych, które przestają, zagregowane krzywe uczenia mogą być
zniekształcone przez ubywanie obserwacji \citep{murray2013,goutte2018}; krzywą obliczono więc
dodatkowo dla zbalansowanej kohorty sekwencji liczących co najmniej \CohortK{} prób. Zależność odsetka
błędów od numeru próby opisano funkcją potęgową \citep{newell1981}. Ponieważ harmonogram powtórek nie
był adaptacyjny, sprawdzono, czy ponowne zadanie pytania zależało od wyniku poprzedniej próby.

\subsection{Modele statystyczne}
Niech $Y_{ijt}$ oznacza poprawność ($1$ albo $0$) odpowiedzi pracownika $i$ na pytanie $j$ w próbie
$t$. Zastosowano trzy modele regresji logistycznej o wspólnej postaci
\begin{equation*}
\operatorname{logit} P(Y_{ijt}=1) = \alpha_i + \gamma_j + \mathbf{x}_{ijt}^{\top}\boldsymbol{\beta},
\end{equation*}
gdzie $\alpha_i$ i $\gamma_j$ są efektami stałymi pracownika i pytania. Efekty stałe pochłaniają
stałe różnice poziomu wiedzy między osobami i trudności między pytaniami, dlatego współczynniki
$\boldsymbol{\beta}$ opisują zmiany w obrębie tej samej osoby i tego samego pytania. Modele różniły
się zbiorem odpowiedzi i wektorem predyktorów $\mathbf{x}_{ijt}$:
\begin{description}
\item[M1, model uczenia (PB1).] Wszystkie odpowiedzi; jedynym predyktorem jest numer próby,
  $\beta_1 \log_2 t$. Wartość $e^{\beta_1}$ jest ilorazem szans (OR) poprawnej odpowiedzi przy
  podwojeniu numeru próby (np.\ z 1 do 2 albo z 4 do 8).
\item[M2, model retencji (PB2).] Wyłącznie próby powtórne ($t \geq 2$);
  $\beta_1 \log_2 t + \beta_2 Y_{ij,t-1} + \beta_3 \log_2 \Delta_{ijt} + \beta_4 S_{it}$, gdzie
  $Y_{ij,t-1}$ to wynik poprzedniej próby, $\Delta_{ijt}$ odstęp od niej w dniach (ograniczony od dołu
  do 1 dnia), a $S_{it}$ staż w latach. Konstrukcją model odpowiada modelom \emph{performance factors
  analysis} i \emph{half-life regression} stosowanym do logów uczących się \citep{pavlik2009,settles2016}.
\item[M3, model nowych pytań.] Wyłącznie pierwsze próby ($t = 1$); jedynym predyktorem jest staż,
  $\beta_4 S_{it}$. Przy efektach stałych pracownika współczynnik porównuje pierwsze kontakty tej samej
  osoby z kolejnymi pytaniami, napotkanymi przy różnym stażu.
\end{description}
Pracowników i pytania, w których wszystkie odpowiedzi były poprawne albo wszystkie błędne, usuwano
iteracyjnie przed estymacją, ponieważ ich efekty stałe nie są identyfikowalne. Modele szacowano metodą
największej wiarygodności, a błędy standardowe klastrowano względem pracowników.

Odporność oszacowania $e^{\beta_1}$ z modelu M1 sprawdzono w analizach wrażliwości: w zbalansowanej
kohorcie, po obcięciu sekwencji do 10 prób, w warstwach długości sekwencji (kontrola obciążenia typu
\emph{pattern-mixture}), po wykluczeniu pracowników z mniej niż 50 odpowiedziami oraz po dodaniu stażu
jako predyktora. Efekt uśredniony populacyjnie oszacowano ponadto metodą uogólnionych równań
estymujących (GEE) z wymienną macierzą korelacji roboczej \citep{liang1986}, a w modelu mieszanym
efekty stałe zastąpiono skrzyżowanymi losowymi wyrazami wolnymi pracowników i pytań \citep{baayen2008}.
Z modelu mieszanego wyznaczono współczynniki korelacji wewnątrzklasowej (ICC) w skali latentnej oraz
brzegowy i warunkowy współczynnik $R^2$ \citep{nakagawa2013}.

\subsection{Przejścia i powtarzane błędy}
Przejścia między kolejnymi próbami tego samego pytania zestawiono w tabeli kontyngencji. Dla dwóch
następujących po sobie błędów porównano odsetek ponownych wyborów tego samego dystraktora z odsetkiem
oczekiwanym przy losowym wyborze spośród zaobserwowanych dystraktorów oraz przy wyborze proporcjonalnym
do popularności każdego dystraktora w danym pytaniu.

Analizy uzupełniające, czyli zmianę poprawności na poziomie pracowników oraz analizę pytań metodami
klasycznej teorii testu, opisano w materiale uzupełniającym. Obliczenia wykonano w języku Python
\VerPython{} z użyciem bibliotek pandas \VerPandas, NumPy \VerNumpy, SciPy \VerScipy, statsmodels
\VerStatsmodels{} i Matplotlib \VerMatplotlib{}.
%  (ziarno generatora liczb losowych \Seed).

\section{Wyniki}

% \subsection{Pierwsza próba}
W pierwszych próbach, czyli przy pierwszym kontakcie pracownika z danym pytaniem, zarejestrowano
\FirstN{} odpowiedzi \NEmployees{} pracowników. Poprawność wynosiła \FirstAcc\% (95\% PU
\FirstAccCI; tab.~\ref{tab:pierwsza_proba_wiek}). W przedziałach wieku mieściła się w zakresie od
\FirstAccMin\% (\FirstAccMinGroup{} lat) do \FirstAccMax\% (\FirstAccMaxGroup{} lat), a przedziały
ufności w grupach o mniejszej liczbie pracowników były szerokie. Po uwzględnieniu trudności pytań
szanse poprawnej odpowiedzi w pierwszej próbie rosły wraz z wiekiem (OR \FirstAgeOR{} na 10 lat,
95\% PU \FirstAgeORCI; $p \FirstAgeP$). Zależność ta nie miała jednak charakteru monotonicznego,
ponieważ najniższą poprawność odnotowano w przedziale \FirstAccMinGroup{} lat, a nie w przedziale
najmłodszym. % SPRAWDŹ: kierunek OR i położenie minimum po zmianie danych

\begin{table}[H]
\centering
\caption{Poprawność w pierwszej próbie ogółem i według przedziału wieku pracownika.}
\label{tab:pierwsza_proba_wiek}
\footnotesize
\setlength{\tabcolsep}{4pt}
\tabelacsv{l*{7}{r}}{pierwsza_proba_wiek}% 7 = ogółem + przedziały wieku + brak danych
\uwagatabeli{Pierwsza próba to pierwsza odpowiedź w sekwencji. 95\% PU wyznaczono metodą bootstrap.
Zależność od wieku oszacowano regresją logistyczną poprawności pierwszej próby względem środka
przedziału wieku (na 10 lat) z efektami stałymi pytania i błędami standardowymi klastrowanymi względem
pracowników (\FirstAgeN{} odpowiedzi \FirstAgeEmp{} pracowników o znanym wieku): OR \FirstAgeOR,
95\% PU \FirstAgeORCI, $p \FirstAgeP$.}
\end{table}

\subsection{Zmiany w sekwencji (PB1)}
Poprawność rosła wraz z numerem próby: z \AccAttOne\% \AccAttOneCI{} w pierwszej próbie do
\AccAttMax\% \AccAttMaxCI{} w próbie \AttMax{}. i kolejnych (rys.~\ref{fig:krzywa_uczenia}).
Liczba sekwencji malała jednak wraz z numerem próby, dlatego zmianę
oceniono również w kohorcie \CohortN{} sekwencji. W kohorcie poprawność wzrosła z \CohortAccOne\%
\CohortAccOneCI{} w próbie~1 do \CohortAccK\% \CohortAccKCI{} w próbie~\CohortK{}, czyli
o \SeqChange{} p.p. (95\% PU \SeqChangeCI; tab.~\ref{tab:sekwencje_wiek}). Wzrost obserwowany w tych
samych sekwencjach wskazuje, że poprawa nie wynika z wcześniejszego kończenia sekwencji z niższą
poprawnością. Odsetek błędów był funkcją potęgową numeru próby (wykładnik \PowerLawB,
$R^2 = \PowerLawRsq$), zgodnie z potęgowym prawem praktyki \citep{newell1981}.

Wzrost poprawności w kohorcie wystąpił we wszystkich przedziałach wieku i wynosił od \SeqChangeMin{}
p.p. (\SeqChangeMinGroup{} lat) do \SeqChangeMax{} p.p. (\SeqChangeMaxGroup{} lat;
tab.~\ref{tab:sekwencje_wiek}). Tempo poprawy nie zależało od wieku: iloraz szans dla interakcji
numeru próby z wiekiem wynosił \SeqAgeOR{} na 10 lat (95\% PU \SeqAgeORCI; $p \SeqAgeP$). Różnice
między przedziałami wieku dotyczyły zatem raczej poziomu wyjściowego niż przyrostu poprawności
w kolejnych próbach. % SPRAWDŹ: znak zmiany we wszystkich przedziałach i istotność interakcji

\begin{figure}[H]
\centering
\includegraphics{\FigDir rys_krzywa_uczenia}
\caption{Poprawność według numeru próby we wszystkich sekwencjach dostępnych w danej próbie (ich liczba
maleje wraz z numerem próby). Słupki błędów: 95\% PU (bootstrap). Linia przerywana: dopasowanie funkcji
potęgowej do odsetka błędów. Próby \MaxAttemptPooled{}. i dalsze połączono. Przebieg w kohorcie
przedstawiono na rys.~\ref{fig:stany_kohorty}.}
\label{fig:krzywa_uczenia}
\end{figure}

\begin{table}[H]
\centering
\caption{Poprawność w kolejnych próbach w kohorcie, ogółem i według przedziału wieku pracownika.}
\label{tab:sekwencje_wiek}
\footnotesize
\setlength{\tabcolsep}{4pt}
\tabelacsv{l*{7}{r}}{sekwencje_wiek}% 7 = ogółem + przedziały wieku + brak danych
\uwagatabeli{Kohorta: sekwencje liczące co najmniej \CohortK{} prób, próby od 1 do \CohortK{}.
p.p.: punkty procentowe; 95\% PU zmiany wyznaczono metodą bootstrap. Zależność tempa poprawy od wieku
oszacowano regresją logistyczną poprawności względem $\log_2(\text{numer próby})$ i jego interakcji
z wiekiem (środek przedziału, na 10 lat), z efektami stałymi pracownika i pytania oraz błędami
standardowymi klastrowanymi względem pracowników (\SeqAgeN{} odpowiedzi pracowników o znanym wieku).}
\end{table}

\subsection{Wynik poprzedniej próby, odstęp i zapominanie (PB2)}
Wzrost poprawności opisany w PB1 jest wynikiem netto dwóch procesów zachodzących między kolejnymi
próbami: korygowania błędów i zapominania odpowiedzi poprawnych. Aby je rozdzielić, każdą próbę od
drugiej wzwyż powiązano z wynikiem próby poprzedniej i odstępem od niej. Harmonogram powtórek nie
zniekształcał tego porównania, ponieważ pytanie wracało po odpowiedzi poprawnej i błędnej z podobnym
prawdopodobieństwem (\ReaskAfterCorrect\% i \ReaskAfterIncorrect\%) i po podobnym odstępie (mediana
\GapAfterCorrect{} i \GapAfterIncorrect{} dni). Odpowiedź poprawna powtórzyła się w kolejnej próbie
w \RetainPct\% \RetainCI{} przypadków, a odpowiedź błędna została skorygowana w \RecoverPct\%
\RecoverCI{} (tab.~\ref{tab:przejscia}). Korekty błędów przeważały zatem nad zapominaniem
(\ForgetPct\%), co tłumaczy przyrost poprawności w kolejnych próbach. Upływ czasu działał
w przeciwnym kierunku, lecz tylko w odniesieniu do odpowiedzi poprawnych: odsetek ich powtórzeń
malał z \RetShortCorrect\% przy odstępie krótszym niż 14 dni do \RetLongCorrect\% przy odstępie
wynoszącym co najmniej 240 dni, podczas gdy odsetek skorygowanych błędów nie wykazywał wyraźnego
trendu (\RetShortIncorrect\% i \RetLongIncorrect\% w tych samych przedziałach; rys.~\ref{fig:odstep}).
Wiedza potwierdzona poprawną odpowiedzią podlegała więc
zapominaniu między powtórkami \citep{cepeda2006,cepeda2008}, natomiast szansa skorygowania błędu po
informacji zwrotnej nie zależała od czasu, który upłynął do kolejnego kontaktu z pytaniem. % SPRAWDŹ: brak trendu dla odpowiedzi błędnych po zmianie danych

\begin{figure}[H]
\centering
\includegraphics{\FigDir rys_odstep}
\caption{Poprawność według odstępu, osobno dla prób poprzedzonych odpowiedzią poprawną (koła)
i błędną (kwadraty). Słupki błędów: 95\% PU (bootstrap).}
\label{fig:odstep}
\end{figure}

\begin{table}[H]
\centering
\caption{Przejścia między kolejnymi próbami tego samego pytania.}
\label{tab:przejscia}
\small
\tabelacsv{lrrlr}{przejscia}
\uwagatabeli{Próby od drugiej wzwyż według wyniku próby poprzedniej. 95\% PU wyznaczono metodą
bootstrap.}
\end{table}

\subsection{Utrwalone błędy (PB3)}
Większość błędów miała charakter przejściowy: w \EverCorrectedPct\% z \FirstWrongSeqN{} sekwencji
rozpoczętych odpowiedzią błędną pracownik odpowiedział poprawnie w którejś z kolejnych prób. Szansa
korekty malała jednak z każdym kolejnym błędem. Po jednym błędzie poprawność wynosiła \CorrRunOne\%
\CorrRunOneCI, po dwóch błędach z rzędu \CorrRunTwo\%, po trzech \CorrRunThree\%, a po co najmniej
\RunMax{} błędach \CorrRunMax\% \CorrRunMaxCI{} (tab.~\ref{tab:seria_bledow}). Każde kolejne
wyświetlenie poprawnej odpowiedzi przynosiło więc mniejszy efekt. Spadek ten może częściowo wynikać
z tego, że wiele błędów z rzędu występowało w trudniejszych pytaniach, nie musi zatem w całości odzwierciedlać
utrwalania się błędu.

O naturze powtórzonych błędów świadczy wybór dystraktora. Gdy pracownik pomylił się w dwóch kolejnych
próbach, w \SameDistractorPct\% \SameDistractorCI{} przypadków ($n = \SameDistractorN$) wskazał ten sam
dystraktor, wobec \SameDistractorExpected\% oczekiwanych przy wyborze losowym
i \SameDistractorExpectedPop\% oczekiwanych przy uwzględnieniu popularności poszczególnych dystraktorów
w danym pytaniu. Powtórzony błąd był zatem wyborem konkretnej odpowiedzi, a nie zgadywaniem. Tendencja
ta nie słabła z upływem czasu: do tego samego dystraktora wracano w \SameDistrGapShort\% przypadków przy
odstępie krótszym niż 30 dni, w \SameDistrGapMid\% przy odstępie od 30 do 89 dni i w \SameDistrGapLong\%
przy odstępie wynoszącym co najmniej 90 dni. Powtórzone błędy skupiały się przy tym w niewielkiej części
materiału, ponieważ na \PersistTopItems{} spośród \PersistItems{} pytań, w których wystąpiły, przypadało
\PersistTopItemsShare\% wszystkich powtórzonych błędów.

Zestawienie wyników PB2 i PB3 ujawnia asymetrię. Odpowiedzi poprawne słabły wraz z odstępem, natomiast
błędne przekonanie, które przetrwało pierwszą korektę, utrzymywało się niezależnie od upływu czasu
i w coraz mniejszym stopniu ulegało informacji zwrotnej. Wyświetlenie poprawnej odpowiedzi wystarczało
do skorygowania większości błędów, ale nie usuwało przyczyny, dla której pracownik wybierał określony
dystraktor. % SPRAWDŹ: monotoniczny spadek korekty i brak trendu dystraktora wg odstępu po zmianie danych

\begin{table}[H]
\centering
\caption{Poprawność według liczby błędów popełnionych z rzędu bezpośrednio przed daną próbą.}
\label{tab:seria_bledow}
\footnotesize
\setlength{\tabcolsep}{5pt}
\tabelacsv{l*{5}{r}}{seria_bledow}% 5 = 0, 1, 2, 3, >=RunMax
\uwagatabeli{Próby od drugiej wzwyż. Kolumna 0: próby poprzedzone odpowiedzią poprawną. 95\% PU
wyznaczono metodą bootstrap.}
\end{table}

\subsection{Podsumowanie wyników}
Wyniki pozwalają odpowiedzieć na trzy pytania badawcze.
\begin{enumerate}[label=(PB\arabic*)]
  \item Poprawność rosła wraz z numerem próby, również w kohorcie, w której zwiększyła się
  o \SeqChange{} p.p. między próbą~1 a próbą~\CohortK{}. Przyrost nie zależał od wieku pracownika,
  a wiek wiązał się jedynie z poprawnością w pierwszej próbie.
  \item Przyrost poprawności wynikał z przewagi korekt błędów (\RecoverPct\%) nad zapominaniem
  (\ForgetPct\%). Zapominanie nasilało się wraz z odstępem, natomiast odsetek korekt błędów od odstępu
  nie zależał.
  \item Większość błędów korygowano (\EverCorrectedPct\% sekwencji rozpoczętych odpowiedzią błędną),
  lecz powtórzone błędy polegały na wyborze tego samego dystraktora (\SameDistractorPct\%), nie słabły
  wraz z odstępem, a szansa ich korekty malała z każdym kolejnym błędem z rzędu, do \CorrRunMax\%
  po co najmniej \RunMax{} błędach.
\end{enumerate}
Łącznie wyniki wskazują, że program utrwalał wiedzę sprawdzaną przez pytania, przy czym wiedza ta
wymagała podtrzymywania kolejnymi powtórkami, a niewielka część pytań generowała błędy odporne na
informację zwrotną ograniczoną do wyświetlenia poprawnej odpowiedzi. % SPRAWDŹ
Rys.~\ref{fig:stany_kohorty} zestawia te wyniki w kohorcie: linia poprawności rośnie (PB1), korekty
błędów przeważają nad zapominaniem (PB2), a powtórzone błędy utrzymują się we wszystkich kolejnych
próbach (PB3).

\begin{figure}[H]
\centering
\includegraphics{\FigDir rys_stany_kohorty}
\caption{Poprawność i stan odpowiedzi w kohorcie (\CohortN{} sekwencji, próby od 1 do \CohortK{}).
Słupki: udział odpowiedzi według wyniku bieżącej i poprzedniej próby (definicje IX--XI
w podrozdziale~\ref{sec:definicje}); w próbie~1 tylko odpowiedzi poprawne i błędne. Linia: poprawność z 95\% PU (bootstrap),
wartości nad słupkami.}
\label{fig:stany_kohorty}
\end{figure}

% \section{Wyniki - stare do usuniecia}

% \subsection{Próba i harmonogram powtórek}
% Zbiór danych obejmował \NResponses{} odpowiedzi \NEmployees{} pracowników na \NItems{} pytań
% w \NModules{} modułach, zebranych od \DateStart{} do \DateEnd{} (tab.~\ref{tab:zbior_danych}).
% Mediana liczby odpowiedzi na pracownika wynosiła \RespPerEmpMedian{} (IQR
% \RespPerEmpQOne--\RespPerEmpQThree), udzielonych na \ItemsPerEmpMedian{} różnych pytań w ciągu
% \ActiveDaysMedian{} dni aktywności, a mediana czasu uczestnictwa w programie wynosiła
% \SpanMonthsMedian{} miesięcy. Okresy bez odpowiedzi odpowiadają zawieszeniom programu na czas
% przeglądu pytań (rys.~\ref{fig:aktywnosc}). Wśród \DemoMatched{} pracowników o znanych danych
% demograficznych kobiety stanowiły \DemoWomenPct\%, a mediana wieku wynosiła \AgeMedian{} lat
% (tab.~\ref{tab:demografia}). Ogólny odsetek poprawnych odpowiedzi wynosił \AccOverall\%, a w pierwszych
% próbach \AccFirst\%.

% Harmonogram powtórek nie zależał od wyniku poprzedniej próby: prawdopodobieństwo kolejnej próby
% wynosiło \ReaskAfterCorrect\% po odpowiedzi poprawnej i \ReaskAfterIncorrect\% po błędnej, natomiast
% mediana odstępu do kolejnej próby wynosiła odpowiednio \GapAfterCorrect{} i \GapAfterIncorrect{} dni.
% Kontynuacja sekwencji nie była więc selekcjonowana według wyników. % SPRAWDŹ

% \subsection{Uczenie się w kolejnych próbach (PB1)}
% Odsetek poprawnych odpowiedzi rósł wraz z numerem próby: z \AccAttOne\% \AccAttOneCI{} w pierwszej
% próbie do \AccAttMax\% \AccAttMaxCI{} w próbie \AttMax{}. i kolejnych (rys.~\ref{fig:krzywa_uczenia},
% tab.~\ref{tab:krzywa_uczenia}). Wzrost był podobny w zbalansowanej kohorcie \CohortN{} sekwencji
% liczących co najmniej \CohortK{} prób (\CohortAccOne\% \CohortAccOneCI{} w próbie~1 wobec \CohortAccK\%
% \CohortAccKCI{} w próbie~\CohortK), nie wynika więc z selektywnej kontynuacji przez pracowników
% osiągających lepsze wyniki. Odsetek błędów był funkcją potęgową numeru próby (wykładnik \PowerLawB,
% $R^2 = \PowerLawRsq$), zgodnie z potęgowym prawem praktyki \citep{newell1981}.

% W modelu uczenia M1 każde podwojenie numeru próby wiązało się z \ORMOneAttempt-krotnie wyższymi
% szansami poprawnej odpowiedzi (95\% PU \ORMOneAttemptCI; tab.~\ref{tab:regresja}). Efekt ten wyjaśniał
% jednak niewielką część zmienności poprawności (brzegowy $R^2 = \RsqMarginal$ wobec warunkowego
% $R^2 = \RsqConditional$), a pytania odpowiadały za większą jej część niż pracownicy (ICC
% \ICCItem{} wobec \ICCEmployee). Na poziomie pracowników wynik w trzeciej próbie był lepszy niż
% w pierwszej u \EmpImproved{} z \EmpN{} osób, a gorszy u \EmpWorsened{} (średnia zmiana
% \EmpMeanChange{} p.p., 95\% PU \EmpMeanChangeCI; tab.~\ref{tab:zmiana_pracownikow}).

% \begin{table*}[htbp]
% \centering
% \caption{Modele regresji logistycznej poprawności odpowiedzi: ilorazy szans [95\% PU].}
% \label{tab:regresja}
% \small
% \setlength{\tabcolsep}{3.5pt}
% \tabelacsv{llll}{regresja}
% \uwagatabeli{Definicje modeli M1--M3 podano w podrozdziale ,,Modele statystyczne''. Wszystkie modele
% zawierają efekty stałe pracownika i pytania; błędy standardowe klastrowane względem pracowników.
% Pracowników i pytania bez zmienności poprawności usunięto przed estymacją (M1: \MOneDropped{}
% odpowiedzi). OR dla numeru próby i odstępu dotyczą podwojenia wartości, dla stażu jednego roku.}
% \end{table*}

% \subsection{Wynik poprzedniej próby, odstęp i zapominanie (PB2)}
% W modelu retencji M2, po uwzględnieniu numeru próby, poprawna poprzednia odpowiedź wiązała się
% z wyższymi szansami poprawnej odpowiedzi (OR \ORMTwoPrev, 95\% PU \ORMTwoPrevCI), natomiast dłuższy
% odstęp od poprzedniej próby wiązał się z niższymi szansami (OR \ORMTwoGap{} na podwojenie odstępu,
% 95\% PU \ORMTwoGapCI; tab.~\ref{tab:regresja}). Wśród odpowiedzi poprzednio poprawnych odsetek
% poprawnych zmniejszał się z \RetShortCorrect\%, gdy pytanie powtórzono w ciągu 14 dni, do
% \RetLongCorrect\% po 240 dniach lub dłużej; odpowiedzi poprzednio błędne były korygowane odpowiednio
% w \RetShortIncorrect\% i \RetLongIncorrect\% przypadków (rys.~\ref{fig:odstep},
% tab.~\ref{tab:retencja_wg_odstepu}). Wzorzec ten jest zgodny z zapominaniem między powtórkami
% \citep{cepeda2006,cepeda2008}. % SPRAWDŹ

% \subsection{Przejścia i utrwalone błędy (PB3)}
% Spośród odpowiedzi poprzednio poprawnych \RetainPct\% \RetainCI{} było poprawnych również w kolejnej
% próbie, a \ForgetPct\% zapomniano. Spośród odpowiedzi poprzednio błędnych \RecoverPct\% \RecoverCI{}
% skorygowano, a \PersistPct\% pozostało błędnych, mimo że po każdej odpowiedzi wyświetlano poprawną
% odpowiedź (tab.~\ref{tab:przejscia}). Gdy błąd się powtarzał (\SameDistractorN{} przypadków),
% pracownicy wybierali ten sam dystraktor w \SameDistractorPct\% \SameDistractorCI{} przypadków, wobec
% \SameDistractorExpected\% oczekiwanych przy wyborze losowym i \SameDistractorExpectedPop\% oczekiwanych
% przy uwzględnieniu popularności poszczególnych dystraktorów. Wskazuje to na utrwalone błędne
% przekonania, a nie na zgadywanie. % SPRAWDŹ

% \subsection{Odporność wyników i nowe pytania}
% Oszacowanie efektu numeru próby pozostawało zbliżone we wszystkich specyfikacjach (OR od \ORRange;
% rys.~\ref{fig:wrazliwosc}, tab.~\ref{tab:wrazliwosc}) i wynosiło \ORSensCohort{} w zbalansowanej
% kohorcie, \ORSensLenThree, \ORSensLenFive{} i \ORSensLenEight{} w sekwencjach liczących 3--4, 5--7
% i co najmniej 8 prób, \ORSensGee{} w uśrednionym populacyjnie modelu GEE oraz \ORSensMixed{} w modelu
% mieszanym. Oszacowanie było wyższe po dodaniu stażu jako predyktora (\ORSensTenure), ponieważ staż
% jest skorelowany z numerem próby.

% W modelu nowych pytań M3 odsetek poprawnych odpowiedzi przy pierwszym kontakcie z pytaniem nie rósł
% wraz ze stażem (OR \ORMThreeTenure{} na rok, 95\% PU \ORMThreeTenureCI; tab.~\ref{tab:regresja}).
% Przedział ufności obejmuje zarówno brak efektu, jak i niewielki spadek, dlatego wynik oznacza brak
% dowodu na uogólnienie poprawy na nowe pytania, a nie dowód jego braku. % SPRAWDŹ

% \begin{figure}[tbp]
% \centering
% \includegraphics{\FigDir rys_wrazliwosc}
% \caption{Iloraz szans (95\% PU) poprawnej odpowiedzi na podwojenie numeru próby w różnych
% specyfikacjach modelu M1. Romb: model główny; linia pionowa: brak efektu.}
% \label{fig:wrazliwosc}
% \end{figure}

% =============================================================================
% Dyskusja i wnioski. Plik edytowany ręcznie -- skrypt analiza_odpowiedzi_naukowa.py go NIE nadpisuje.
% Liczby: definicje.tex; literatura: references.bib.
% =============================================================================
\section{Dyskusja}

\subsection{Główne wyniki}
W ciągu trzech lat rutynowego stosowania programu mikrouczenia w miejscu pracy odsetek poprawnych
odpowiedzi na wielokrotnie prezentowane pytania rósł systematycznie wraz z liczbą przypomnień, zgodnie
z funkcją potęgową, a efekt był odporny na ubywanie obserwacji, specyfikację modelu i klastrowanie
danych. Poprawność malała wraz z odstępem od poprzedniego przypomnienia, wcześniejsze błędy były
w kolejnej próbie korygowane tylko częściowo, a powtarzane błędy w większości odtwarzały ten sam
dystraktor. Poprawność odpowiedzi na nowe pytania nie rosła wraz ze stażem. % SPRAWDŹ

\subsection{Porównanie z wcześniejszymi badaniami}
Stopniowy, regularny przyrost przypomina krzywe uczenia obserwowane w wielu zbiorach danych
edukacyjnych \citep{koedinger2023} i jest zgodny z efektami testowania i rozłożenia w czasie
\citep{roediger2006,cepeda2006}. Randomizowane badania edukacji rozłożonej w czasie wykazały lepsze
utrwalenie wiedzy \citep{kerfoot2007mededuc,kerfoot2009}; nasze dane obserwacyjne rozszerzają te
wyniki na niemedyczne środowisko pracy z długimi odstępami między powtórkami. Ujemny związek między
odstępem a poprawnością jest zgodny z modelami krzywej zapominania dopasowywanymi do logów uczących się
\citep{settles2016,tabibian2019}. Podobnie jak w badaniach, w których wielokrotnie wykorzystywano te
same pytania \citep{donker2022}, poprawy w powtarzanych pytaniach nie da się oddzielić od zapamiętania
konkretnego pytania i jego odpowiedzi; brak efektu stażu w nowych pytaniach sugeruje ograniczony
transfer.

\subsection{Implikacje praktyczne}
Program ponownie prezentował pytania w tym samym odstępie niezależnie od wyniku poprzedniej próby.
Harmonogramy adaptacyjne, które skracają odstęp po błędach i wydłużają go po poprawnych odpowiedziach,
efektywniej wykorzystują historię uczącego się \citep{lindsey2014,settles2016,tabibian2019}. Konkretne
dystraktory utrzymywały się, mimo że po każdej odpowiedzi wyświetlano wyjaśnienie poprawnej odpowiedzi,
co sugeruje, że ogólne wyjaśnienie nie wystarcza do skorygowania błędnych przekonań; skuteczniejsza może
być informacja zwrotna odnosząca się do konkretnego wybranego dystraktora. Analiza pytań wskazała
kandydatów do zmiany: pytania o niskiej dyskryminacji, pytania, na które niemal wszyscy odpowiadali
poprawnie już w pierwszej próbie (mało miejsca na uczenie się), oraz dystraktory niefunkcjonalne
\citep{tarrant2009,gierl2017}.

\subsection{Ograniczenia}
Badanie ma charakter obserwacyjny i nie obejmuje grupy kontrolnej, nie można więc przypisać przyrostu
wiedzy programowi w sensie przyczynowym; poprawa w powtarzanych pytaniach może częściowo odzwierciedlać
rozpoznawanie pytania, a nie jego zrozumienie. Udział i aktywność różniły się między pracownikami,
a wyniki dla sekwencji z wieloma próbami opisują bardziej aktywnych uczestników, choć analizy
zbalansowanej kohorty i długości sekwencji nie wskazały na obciążenie wynikające z ubywania obserwacji.
Wysoka poprawność w pierwszych próbach (\AccFirst\%) ogranicza możliwość poprawy (efekt sufitowy).
Dystraktory, których nigdy nie wybrano, nie są widoczne w logu, więc funkcjonowanie dystraktorów może
być przeszacowane. Nie dysponowano informacjami o stanowiskach pracowników, ich wcześniejszej wiedzy
ani danych demograficznych. Znaczniki czasu nie zawierały informacji o strefie czasowej. Model mieszany
zawierał wyłącznie losowe wyrazy wolne; nie modelowano indywidualnego tempa uczenia się pracowników.
Dane pochodzą z jednej organizacji. \fillin{Dodaj ograniczenia specyficzne dla kontekstu.}

\section{Wnioski}
Rutynowemu, powtarzanemu przypominaniu w szkoleniach pracowniczych towarzyszyła systematyczna poprawa
specyficzna dla powtarzanych pytań oraz zapominanie między powtórkami. Dane z logów takich programów
mogą służyć do monitorowania uczenia się, oceny pytań i projektowania adaptacyjnych harmonogramów
powtórek.


\section*{Oświadczenia}

% SPRAWDŹ: sposób kodowania danych pracowników oraz status zgody
% komisji etycznej lub zwolnienia z obowiązku jej uzyskania.
\noindent\textbf{Etyka.}
W badaniu wykorzystano dane przekazane przez bank, w których pracownicy
byli reprezentowani przez kody. Dane dotyczące wieku i płci raportowano
wyłącznie w postaci zagregowanej.
\fillin{Informacja o zgodzie właściwej komisji etycznej lub zwolnieniu
z obowiązku jej uzyskania, jeśli jest wymagana.}\\
\textbf{Dostępność danych.}
Dane wykorzystane w badaniu stanowią własność klienta i nie są
udostępniane publicznie.\\
\textbf{Finansowanie.}
Badanie nie otrzymało finansowania.\\
\textbf{Konflikt interesów.}
Autorzy deklarują brak konfliktu interesów.\\
\textbf{Wykorzystanie narzędzi AI.}
Narzędzia generatywnej sztucznej inteligencji wykorzystano do
przygotowania opisu wyników analizy oraz tłumaczenia tekstu
na język angielski.
\bibliography{references}
\end{document}
